\chapter{The Levi--Civita Connection}
\label{ch:vii22-levi-civita}

A Riemannian metric measures lengths and angles inside each tangent space.  To
differentiate vector fields one also needs a connection.  The fundamental theorem
of Riemannian geometry says that the metric selects a unique connection once two
natural requirements are imposed:
\[
\boxed{\nabla g=0,\qquad T=0.}
\]
This distinguished connection is the Levi--Civita connection.

\section{Metric compatibility}

\begin{definition}[Metric-compatible connection]\label{def:vii22-metric-compatible}
Let \((M,g)\) be Riemannian.  A connection \(\nabla\) on \(TM\) is metric compatible
if
\[
X\,g(Y,Z)
=
g(\nabla_XY,Z)+g(Y,\nabla_XZ)
\]
for all vector fields \(X,Y,Z\).
\end{definition}

\begin{proposition}[Preservation of norms]\label{prop:vii22-norm-derivative}
For a vector field \(V\) along a curve \(\gamma\),
\[
\frac{d}{dt}g(V,V)
=
2g\!\left(\frac{DV}{dt},V\right).
\]
In particular, a parallel vector field has constant length.
\end{proposition}

\begin{proof}
Apply metric compatibility with \(X=\dot\gamma\) and \(Y=Z=V\).
\end{proof}

\section{Torsion-free connections}

Recall from VII/21 that
\[
T(X,Y)=\nabla_XY-\nabla_YX-[X,Y].
\]

\begin{definition}[Torsion free]\label{def:vii22-torsion-free}
A connection is torsion free if \(T=0\), equivalently
\[
\nabla_XY-\nabla_YX=[X,Y].
\]
\end{definition}

In a coordinate frame, where \([\partial_i,\partial_j]=0\), torsion freeness is
equivalent to
\[
\Gamma^k_{ij}=\Gamma^k_{ji}.
\]

\section{The fundamental theorem}

\begin{theorem}[Fundamental theorem of Riemannian geometry]\label{thm:vii22-fundamental}
Every Riemannian manifold \((M,g)\) admits a unique affine connection \(\nabla\)
which is both metric compatible and torsion free.
\end{theorem}

The key formula is obtained by cycling metric compatibility and using zero torsion.

\begin{theorem}[Koszul formula]\label{thm:vii22-koszul}
The Levi--Civita connection satisfies
\[
\begin{aligned}
2g(\nabla_XY,Z)
={}&Xg(Y,Z)+Yg(Z,X)-Zg(X,Y)\\
&-g(X,[Y,Z])+g(Y,[Z,X])+g(Z,[X,Y]).
\end{aligned}
\]
Conversely this identity defines the Levi--Civita connection.
\end{theorem}

\begin{proof}
Write metric compatibility for the cyclic triples
\((X,Y,Z)\), \((Y,Z,X)\), and \((Z,X,Y)\). Add the first two equations and subtract
the third. Replace expressions such as \(\nabla_XY-\nabla_YX\) by \([X,Y]\).
This gives the formula.

For uniqueness, the right side depends only on \(g,X,Y,Z\), so it determines
\(g(\nabla_XY,Z)\) for every \(Z\), hence determines \(\nabla_XY\).  Conversely,
the formula can be checked to define a connection satisfying both required
properties.
\end{proof}

\section{Christoffel symbols from the metric}

Let
\[
g_{ij}=g(\partial_i,\partial_j),
\qquad
(g^{ij})=(g_{ij})^{-1}.
\]

\begin{theorem}[Christoffel formula]\label{thm:vii22-christoffel}
For the Levi--Civita connection,
\[
\Gamma^k_{ij}
=
\frac12 g^{k\ell}
\left(
\partial_i g_{j\ell}
+\partial_j g_{i\ell}
-\partial_\ell g_{ij}
\right).
\]
\end{theorem}

\begin{proof}
Apply the Koszul formula to the commuting coordinate vector fields
\(\partial_i,\partial_j,\partial_\ell\):
\[
2g_{\ell k}\Gamma^k_{ij}
=
\partial_i g_{j\ell}
+\partial_j g_{i\ell}
-\partial_\ell g_{ij}.
\]
Multiply by \(g^{m\ell}\).
\end{proof}

\begin{corollary}[Symmetry]\label{cor:vii22-christoffel-symmetry}
\[
\Gamma^k_{ij}=\Gamma^k_{ji}.
\]
\end{corollary}

\section{Euclidean space}

\begin{example}[Cartesian coordinates]\label{ex:vii22-cartesian}
For the Euclidean metric
\[
g=\sum_i (dx^i)^2,
\]
all \(g_{ij}\) are constant, so every Christoffel symbol vanishes:
\[
\Gamma^k_{ij}=0.
\]
\end{example}

\begin{warning}[Vanishing coefficients are coordinate dependent]\label{warn:vii22-coordinates}
The Euclidean connection has nonzero Christoffel symbols in curvilinear
coordinates.  A connection is geometric; its coefficient array is not a tensor.
\end{warning}

\begin{example}[Polar coordinates]\label{ex:vii22-polar}
For
\[
g=dr^2+r^2d\theta^2,
\]
the nonzero Christoffel symbols are
\[
\Gamma^r_{\theta\theta}=-r,
\qquad
\Gamma^\theta_{r\theta}
=
\Gamma^\theta_{\theta r}
=
\frac1r.
\]
\end{example}

\section{Conformal metrics}

\begin{example}[Conformal Euclidean metric]\label{ex:vii22-conformal}
On a domain in \(\mathbb R^n\), let
\[
g_{ij}=e^{2\varphi}\delta_{ij}.
\]
Then
\[
\Gamma^k_{ij}
=
\delta^k_j\,\partial_i\varphi
+\delta^k_i\,\partial_j\varphi
-\delta_{ij}\,\partial^k\varphi.
\]
\end{example}

\section{Tangential projection for submanifolds}

\begin{theorem}[Induced Levi--Civita connection]\label{thm:vii22-projection}
Let \(M\subset\mathbb R^N\) be an embedded submanifold with induced metric.
For tangent vector fields \(X,Y\),
\[
\nabla_XY=(D_XY)^\top,
\]
the tangential projection of the ambient derivative.
\end{theorem}

\begin{proof}
The projected derivative is a connection.  Euclidean differentiation preserves the
ambient inner product, so its tangential part is metric compatible.  Since the
ambient derivative has zero torsion and the Lie bracket of tangent fields is
tangent, the projected connection is torsion free.  Uniqueness identifies it with
the Levi--Civita connection.
\end{proof}

For a surface in \(\mathbb R^3\), the complementary normal part gives the Gauss
formula
\[
D_XY=\nabla_XY+II(X,Y)N.
\]

\section{Covariant Hessian}

\begin{definition}[Hessian]\label{def:vii22-hessian}
For \(f\in C^\infty(M)\),
\[
\operatorname{Hess}f(X,Y)
=
g(\nabla_X\operatorname{grad}f,Y)
=
X(Yf)-(\nabla_XY)f.
\]
\end{definition}

\begin{proposition}[Symmetry of the Hessian]\label{prop:vii22-hessian-symmetric}
For the Levi--Civita connection,
\[
\operatorname{Hess}f(X,Y)=\operatorname{Hess}f(Y,X).
\]
\end{proposition}

\begin{proof}
Subtract the two expressions:
\[
X(Yf)-Y(Xf)-(\nabla_XY-\nabla_YX)f.
\]
Torsion freeness turns the second parenthesis into \([X,Y]\), and the result
vanishes.
\end{proof}

\section{Why uniqueness matters}

The Levi--Civita connection removes the arbitrary choice present in VII/21.
Once \(g\) is fixed, covariant differentiation, geodesics, parallel transport, and
the Riemann curvature tensor become canonical Riemannian constructions.

\section{Exercises}

\begin{exercise}\label{exr:vii22-01}
State metric compatibility.
\end{exercise}
\begin{hint}
% PEDAGOGY-ENRICHED-VII
Differentiate an inner product. Make the controlling geometric criterion explicit before the calculation. Use metric compatibility together with zero torsion to determine the Levi-Civita connection; in coordinates this becomes the Christoffel-symbol formula.
\end{hint}
\begin{solution}
It is \(Xg(Y,Z)=g(\nabla_XY,Z)+g(Y,\nabla_XZ)\).
\end{solution}


\begin{exercise}\label{exr:vii22-02}
State torsion freeness.
\end{exercise}
\begin{hint}
% PEDAGOGY-ENRICHED-VII
Use the Lie bracket. Work in a convenient local model first, then return to the invariant statement. Use metric compatibility together with zero torsion to determine the Levi-Civita connection; in coordinates this becomes the Christoffel-symbol formula.
\end{hint}
\begin{solution}
It is \(\nabla_XY-\nabla_YX=[X,Y]\).
\end{solution}


\begin{exercise}\label{exr:vii22-03}
What characterizes the Levi--Civita connection?
\end{exercise}
\begin{hint}
% PEDAGOGY-ENRICHED-VII
Combine two conditions. After the main computation, check the hypotheses and geometric interpretation. Use metric compatibility together with zero torsion to determine the Levi-Civita connection; in coordinates this becomes the Christoffel-symbol formula.
\end{hint}
\begin{solution}
It is the unique connection that is metric compatible and torsion free.
\end{solution}


\begin{exercise}\label{exr:vii22-04}
Show a parallel field has constant norm for a metric-compatible connection.
\end{exercise}
\begin{hint}
% PEDAGOGY-ENRICHED-VII
Differentiate \(g(V,V)\). Make the controlling geometric criterion explicit before the calculation. Use metric compatibility together with zero torsion to determine the Levi-Civita connection; in coordinates this becomes the Christoffel-symbol formula.
\end{hint}
\begin{solution}
Metric compatibility gives \(\frac d{dt}g(V,V)=2g(DV/dt,V)=0\).
\end{solution}


\begin{exercise}\label{exr:vii22-05}
Show two parallel fields have constant inner product.
\end{exercise}
\begin{hint}
% PEDAGOGY-ENRICHED-VII
Differentiate \(g(V,W)\). Work in a convenient local model first, then return to the invariant statement. Use metric compatibility together with zero torsion to determine the Levi-Civita connection; in coordinates this becomes the Christoffel-symbol formula.
\end{hint}
\begin{solution}
Both covariant derivative terms vanish, so the inner product is constant.
\end{solution}


\begin{exercise}\label{exr:vii22-06}
Write the Koszul formula.
\end{exercise}
\begin{hint}
% PEDAGOGY-ENRICHED-VII
Cycle metric compatibility. After the main computation, check the hypotheses and geometric interpretation. Use metric compatibility together with zero torsion to determine the Levi-Civita connection; in coordinates this becomes the Christoffel-symbol formula.
\end{hint}
\begin{solution}
The formula expresses \(2g(\nabla_XY,Z)\) as three metric derivatives plus three Lie-bracket terms.
\end{solution}


\begin{exercise}\label{exr:vii22-07}
Why does the Koszul formula prove uniqueness?
\end{exercise}
\begin{hint}
% PEDAGOGY-ENRICHED-VII
Vary \(Z\). Make the controlling geometric criterion explicit before the calculation. Use metric compatibility together with zero torsion to determine the Levi-Civita connection; in coordinates this becomes the Christoffel-symbol formula.
\end{hint}
\begin{solution}
It determines the inner product of \(\nabla_XY\) with every \(Z\), hence determines the vector.
\end{solution}


\begin{exercise}\label{exr:vii22-08}
Why do coordinate Lie-bracket terms vanish in the Koszul formula?
\end{exercise}
\begin{hint}
% PEDAGOGY-ENRICHED-VII
Coordinate vector fields commute. Work in a convenient local model first, then return to the invariant statement. Use metric compatibility together with zero torsion to determine the Levi-Civita connection; in coordinates this becomes the Christoffel-symbol formula.
\end{hint}
\begin{solution}
\([\partial_i,\partial_j]=0\).
\end{solution}


\begin{exercise}\label{exr:vii22-09}
State the Christoffel formula for the Levi--Civita connection.
\end{exercise}
\begin{hint}
% PEDAGOGY-ENRICHED-VII
Use \(g^{k\ell}\). After the main computation, check the hypotheses and geometric interpretation. Use metric compatibility together with zero torsion to determine the Levi-Civita connection; in coordinates this becomes the Christoffel-symbol formula.
\end{hint}
\begin{solution}
\(\Gamma^k_{ij}=\frac12g^{k\ell}(\partial_i g_{j\ell}+\partial_jg_{i\ell}-\partial_\ell g_{ij})\).
\end{solution}


\begin{exercise}\label{exr:vii22-10}
Show the Levi--Civita Christoffel symbols are symmetric in \(i,j\).
\end{exercise}
\begin{hint}
% PEDAGOGY-ENRICHED-VII
Inspect the formula. Make the controlling geometric criterion explicit before the calculation. Use metric compatibility together with zero torsion to determine the Levi-Civita connection; in coordinates this becomes the Christoffel-symbol formula.
\end{hint}
\begin{solution}
The right side is unchanged when \(i\) and \(j\) are interchanged.
\end{solution}


\begin{exercise}\label{exr:vii22-11}
Compute the Christoffel symbols of Euclidean Cartesian coordinates.
\end{exercise}
\begin{hint}
% PEDAGOGY-ENRICHED-VII
The metric coefficients are constant. Work in a convenient local model first, then return to the invariant statement. Use metric compatibility together with zero torsion to determine the Levi-Civita connection; in coordinates this becomes the Christoffel-symbol formula.
\end{hint}
\begin{solution}
All Christoffel symbols vanish.
\end{solution}


\begin{exercise}\label{exr:vii22-12}
Why can Euclidean Christoffel symbols be nonzero in polar coordinates?
\end{exercise}
\begin{hint}
% PEDAGOGY-ENRICHED-VII
Connection coefficients are not tensor components. After the main computation, check the hypotheses and geometric interpretation. Use metric compatibility together with zero torsion to determine the Levi-Civita connection; in coordinates this becomes the Christoffel-symbol formula.
\end{hint}
\begin{solution}
The coordinate frame itself varies, producing nonzero coefficients even for the flat Euclidean connection.
\end{solution}


\begin{exercise}\label{exr:vii22-13}
Compute \(\Gamma^r_{\theta\theta}\) for \(dr^2+r^2d\theta^2\).
\end{exercise}
\begin{hint}
% PEDAGOGY-ENRICHED-VII
Differentiate \(g_{\theta\theta}=r^2\). Make the controlling geometric criterion explicit before the calculation. Use metric compatibility together with zero torsion to determine the Levi-Civita connection; in coordinates this becomes the Christoffel-symbol formula.
\end{hint}
\begin{solution}
\(\Gamma^r_{\theta\theta}=-\frac12\partial_r(r^2)=-r\).
\end{solution}


\begin{exercise}\label{exr:vii22-14}
Compute \(\Gamma^\theta_{r\theta}\) in polar coordinates.
\end{exercise}
\begin{hint}
% PEDAGOGY-ENRICHED-VII
Use \(g^{\theta\theta}=1/r^2\). Work in a convenient local model first, then return to the invariant statement. Use metric compatibility together with zero torsion to determine the Levi-Civita connection; in coordinates this becomes the Christoffel-symbol formula.
\end{hint}
\begin{solution}
\(\Gamma^\theta_{r\theta}=\frac12r^{-2}\partial_r(r^2)=1/r\).
\end{solution}


\begin{exercise}\label{exr:vii22-15}
Show the projected ambient derivative defines a connection on an embedded submanifold.
\end{exercise}
\begin{hint}
% PEDAGOGY-ENRICHED-VII
Project after applying \(D_X\). After the main computation, check the hypotheses and geometric interpretation. Use metric compatibility together with zero torsion to determine the Levi-Civita connection; in coordinates this becomes the Christoffel-symbol formula.
\end{hint}
\begin{solution}
Linearity in the direction and the Leibniz rule survive orthogonal projection onto the tangent bundle.
\end{solution}


\begin{exercise}\label{exr:vii22-16}
Why is the projected connection metric compatible?
\end{exercise}
\begin{hint}
% PEDAGOGY-ENRICHED-VII
Use the Euclidean product rule. Make the controlling geometric criterion explicit before the calculation. Use metric compatibility together with zero torsion to determine the Levi-Civita connection; in coordinates this becomes the Christoffel-symbol formula.
\end{hint}
\begin{solution}
The normal parts are orthogonal to tangent vectors, leaving exactly the intrinsic metric-compatibility identity.
\end{solution}


\begin{exercise}\label{exr:vii22-17}
Why is the projected connection torsion free?
\end{exercise}
\begin{hint}
% PEDAGOGY-ENRICHED-VII
Compare ambient derivatives. Work in a convenient local model first, then return to the invariant statement. Use metric compatibility together with zero torsion to determine the Levi-Civita connection; in coordinates this becomes the Christoffel-symbol formula.
\end{hint}
\begin{solution}
\((D_XY-D_YX)^\top=[X,Y]\), since the bracket is tangent.
\end{solution}


\begin{exercise}\label{exr:vii22-18}
State the Gauss formula for a surface.
\end{exercise}
\begin{hint}
% PEDAGOGY-ENRICHED-VII
Split ambient derivative into tangent and normal parts. After the main computation, check the hypotheses and geometric interpretation. Use metric compatibility together with zero torsion to determine the Levi-Civita connection; in coordinates this becomes the Christoffel-symbol formula.
\end{hint}
\begin{solution}
\(D_XY=\nabla_XY+II(X,Y)N\).
\end{solution}


\begin{exercise}\label{exr:vii22-19}
Define the Hessian of a function.
\end{exercise}
\begin{hint}
% PEDAGOGY-ENRICHED-VII
Differentiate the gradient covariantly. Make the controlling geometric criterion explicit before the calculation. Use metric compatibility together with zero torsion to determine the Levi-Civita connection; in coordinates this becomes the Christoffel-symbol formula.
\end{hint}
\begin{solution}
\(\operatorname{Hess}f(X,Y)=g(\nabla_X\operatorname{grad}f,Y)\).
\end{solution}


\begin{exercise}\label{exr:vii22-20}
Give the equivalent Hessian formula without the gradient.
\end{exercise}
\begin{hint}
% PEDAGOGY-ENRICHED-VII
Expand using metric compatibility. Work in a convenient local model first, then return to the invariant statement. Use metric compatibility together with zero torsion to determine the Levi-Civita connection; in coordinates this becomes the Christoffel-symbol formula.
\end{hint}
\begin{solution}
\(\operatorname{Hess}f(X,Y)=X(Yf)-(\nabla_XY)f\).
\end{solution}


\begin{exercise}\label{exr:vii22-21}
Why is the Levi--Civita Hessian symmetric?
\end{exercise}
\begin{hint}
% PEDAGOGY-ENRICHED-VII
Subtract the two orders. After the main computation, check the hypotheses and geometric interpretation. Use metric compatibility together with zero torsion to determine the Levi-Civita connection; in coordinates this becomes the Christoffel-symbol formula.
\end{hint}
\begin{solution}
The difference equals \([X,Y]f-(\nabla_XY-\nabla_YX)f=0\) by torsion freeness.
\end{solution}


\begin{exercise}\label{exr:vii22-22}
For constant metric coefficients, what does the Christoffel formula give?
\end{exercise}
\begin{hint}
% PEDAGOGY-ENRICHED-VII
All first derivatives vanish. Make the controlling geometric criterion explicit before the calculation. Use metric compatibility together with zero torsion to determine the Levi-Civita connection; in coordinates this becomes the Christoffel-symbol formula.
\end{hint}
\begin{solution}
It gives \(\Gamma^k_{ij}=0\) in that coordinate system.
\end{solution}


\begin{exercise}\label{exr:vii22-23}
What data of VII/21 become canonical after choosing \(g\)?
\end{exercise}
\begin{hint}
% PEDAGOGY-ENRICHED-VII
Use uniqueness. Work in a convenient local model first, then return to the invariant statement. Use metric compatibility together with zero torsion to determine the Levi-Civita connection; in coordinates this becomes the Christoffel-symbol formula.
\end{hint}
\begin{solution}
The general connection is replaced by the unique Levi--Civita connection, making its covariant derivative canonical.
\end{solution}


\begin{exercise}\label{exr:vii22-24}
State the bridge to geodesics.
\end{exercise}
\begin{hint}
% PEDAGOGY-ENRICHED-VII
Use covariant acceleration. After the main computation, check the hypotheses and geometric interpretation. Use metric compatibility together with zero torsion to determine the Levi-Civita connection; in coordinates this becomes the Christoffel-symbol formula.
\end{hint}
\begin{solution}
VII/23 defines geodesics by the Levi--Civita equation \(\nabla_{\dot\gamma}\dot\gamma=0\).
\end{solution}

\section{Solved problem dossiers}

\begin{problem}[Fundamental-theorem problem]\label{prob:vii22-01}
Derive the Koszul formula from metric compatibility and torsion freeness.
\end{problem}
\begin{solution}
Write the three cyclic metric-compatibility identities, add the first two and subtract the third, then replace antisymmetric covariant-derivative differences with Lie brackets.
\end{solution}


\begin{problem}[Uniqueness problem]\label{prob:vii22-02}
Prove at most one metric-compatible torsion-free connection can exist.
\end{problem}
\begin{solution}
Both must satisfy the same Koszul formula, which determines \(g(\nabla_XY,Z)\) for every \(Z\). Nondegeneracy of \(g\) gives equality.
\end{solution}


\begin{problem}[Existence problem]\label{prob:vii22-03}
Explain why the Koszul formula defines a connection.
\end{problem}
\begin{solution}
For fixed \(X,Y\), its right side is \(C^\infty\)-linear in \(Z\), hence represents a unique vector. Checking the connection axioms, metric compatibility, and zero torsion is direct.
\end{solution}


\begin{problem}[Christoffel problem]\label{prob:vii22-04}
Derive the coordinate Christoffel formula from Koszul.
\end{problem}
\begin{solution}
Set \(X=\partial_i,Y=\partial_j,Z=\partial_\ell\), use commuting coordinate fields, then raise the \(\ell\) index with \(g^{k\ell}\).
\end{solution}


\begin{problem}[Polar-coordinate problem]\label{prob:vii22-05}
Compute all nonzero Euclidean Christoffel symbols in polar coordinates.
\end{problem}
\begin{solution}
For \(g_{rr}=1,g_{\theta\theta}=r^2\), only the derivative \(\partial_r g_{\theta\theta}=2r\) is nonzero, yielding \(\Gamma^r_{\theta\theta}=-r\) and \(\Gamma^\theta_{r\theta}=\Gamma^\theta_{\theta r}=1/r\).
\end{solution}


\begin{problem}[Conformal-metric problem]\label{prob:vii22-06}
Derive the Christoffel symbols for \(g=e^{2\varphi}\delta\).
\end{problem}
\begin{solution}
Substitute \(g_{ij}=e^{2\varphi}\delta_{ij}\) and \(g^{ij}=e^{-2\varphi}\delta^{ij}\) into the metric formula and cancel the common factor.
\end{solution}


\begin{problem}[Submanifold problem]\label{prob:vii22-07}
Show tangential projection of the Euclidean derivative is the induced Levi--Civita connection.
\end{problem}
\begin{solution}
Verify connection axioms, metric compatibility, and zero torsion; uniqueness then completes the argument.
\end{solution}


\begin{problem}[Gauss-formula problem]\label{prob:vii22-08}
Recover the second fundamental form from the ambient/intrinsic derivative split.
\end{problem}
\begin{solution}
The normal component of \(D_XY\) is \(\langle D_XY,N\rangle N=II(X,Y)N\), leaving the Levi--Civita tangential component.
\end{solution}


\begin{problem}[Hessian-symmetry problem]\label{prob:vii22-09}
Prove the Hessian is symmetric precisely because the Levi--Civita torsion vanishes.
\end{problem}
\begin{solution}
Compute \(\operatorname{Hess}f(X,Y)-\operatorname{Hess}f(Y,X)\); the result is \([X,Y]f-(\nabla_XY-\nabla_YX)f=-T(X,Y)f\).
\end{solution}


\begin{problem}[Metric-compatibility problem]\label{prob:vii22-10}
Show parallel transport for the Levi--Civita connection must preserve inner products.
\end{problem}
\begin{solution}
Along a curve, differentiate \(g(V,W)\). Metric compatibility expresses the derivative as the sum of pairings with \(DV/dt\) and \(DW/dt\), both zero for parallel fields.
\end{solution}


\begin{problem}[Normal-coordinate preview]\label{prob:vii22-11}
Why is it plausible that Levi--Civita Christoffel symbols can be made zero at one point?
\end{problem}
\begin{solution}
The connection coefficients are not tensors and can be canceled at a chosen point by a suitable coordinate choice. VII/23 constructs such coordinates via the exponential map.
\end{solution}


\begin{problem}[Connection-versus-curvature problem]\label{prob:vii22-12}
Why does zero Christoffel symbol at one point not imply zero curvature?
\end{problem}
\begin{solution}
Curvature depends on first derivatives and quadratic combinations of connection coefficients. Coordinates can kill the coefficients at one point but cannot generally kill their derivatives there.
\end{solution}

\section{Challenge problems}

\begin{problem}[Difference of metric connections]\label{prob:vii22-ch01}
Suppose \(\nabla,\widetilde\nabla\) are both metric compatible. Describe the algebraic constraint on \(A(X)Y=\widetilde\nabla_XY-\nabla_XY\).
\end{problem}
\begin{solution}
Metric compatibility for both connections implies \(g(A(X)Y,Z)+g(Y,A(X)Z)=0\). Thus for each \(X\), the endomorphism \(A(X)\) is skew-adjoint.
\end{solution}


\begin{problem}[Uniqueness via difference tensor]\label{prob:vii22-ch02}
Use the previous fact plus zero torsion to prove uniqueness without invoking the full Koszul formula.
\end{problem}
\begin{solution}
Torsion freeness gives \(A(X)Y=A(Y)X\), while metric compatibility gives skew-adjointness in the last two slots. Cycling the three arguments forces the trilinear form \(g(A(X)Y,Z)\) to equal its own negative, hence vanish.
\end{solution}


\begin{problem}[Conformal connection difference]\label{prob:vii22-ch03}
Compare the Levi--Civita connections of \(g\) and \(\widetilde g=e^{2\varphi}g\).
\end{problem}
\begin{solution}
The difference is \(\widetilde\nabla_XY-\nabla_XY=X(\varphi)Y+Y(\varphi)X-g(X,Y)\operatorname{grad}_g\varphi\). This follows from the Koszul formula or the conformal Christoffel formula.
\end{solution}


\begin{problem}[Normal coordinates preview]\label{prob:vii22-ch04}
Assume coordinates satisfy \(g_{ij}(p)=\delta_{ij}\) and \(\partial_k g_{ij}(p)=0\). What are the Levi--Civita coefficients at \(p\)?
\end{problem}
\begin{solution}
The Christoffel formula is built from first derivatives of \(g_{ij}\), so every \(\Gamma^k_{ij}(p)\) vanishes.
\end{solution}


\begin{problem}[Variational bridge]\label{prob:vii22-ch05}
Explain why metric compatibility will force a geodesic to have constant speed.
\end{problem}
\begin{solution}
For a geodesic, \(D\dot\gamma/dt=0\). Metric compatibility then gives \(\frac d{dt}g(\dot\gamma,\dot\gamma)=2g(D\dot\gamma/dt,\dot\gamma)=0\).
\end{solution}

\section*{Bridge to VII/23}
With the Levi--Civita connection fixed canonically by the metric, VII/23 studies its autoparallel curves: geodesics, the exponential map, normal coordinates, and variational length geometry.
